\section{بررسی درستی طراحی}
\label{sec:book-verification}
\subsection{رابطه میان \lr{Specification} و نمایش‌های طراحی}
\lr{Verification} بررسی می‌کند که طراحی با رفتار و محدودیت‌های مورد انتظار
سازگار باشد. \lr{Specification} باید رفتار خروجی، شرایط ورودی، \lr{Reset}،
زمان‌بندی و سایر قیود لازم را مشخص کند. با دقیق‌ترشدن نمایش، تعداد جزئیات و هزینه
تحلیل افزایش می‌یابد. کتاب برای مدیریت این هزینه، بررسی عملکرد در سطوح بالاتر و
بررسی سازگاری میان نمایش‌های متوالی را توضیح می‌دهد
\cite[بخش \secnum{۱.۱۱}، شکل \secnum{۱.۷۰}]{weste}.

در سطح معماری، مدل رفتاری مانند برنامه \lr{C} با تعریف معماری سنجیده می‌شود.
در سطح منطق، \lr{HDL} با ورودی و خروجی مورد انتظار بررسی می‌شود. در مرحله مدار،
عملکرد شبکه ترانزیستور با تابع منطقی مقایسه می‌شود. در مرحله فیزیکی، اتصال‌های
استخراج‌شده از \lr{Layout} با \lr{Schematic} مقایسه می‌شوند. در کنار این زنجیره،
مدار و \lr{Layout} باید قیود زمان‌بندی و توان را برآورده کنند.

برای نمونه، اگر نیاز سیستم جمع هشت‌بیتی با \lr{Carry-out} باشد، رفتار مرجع باید
جمع کامل ۹ بیتی را مشخص کند. سپس خروجی هشت‌بیتی، هشت بیت کم‌ارزش نتیجه و
\lr{Carry-out} بیت نهم است. اگر تعریف علامت‌دار بودن داده یا رفتار \lr{Overflow}
مبهم بماند، پاسخ مورد انتظار آزمون نیز مبهم می‌شود.

\subsection{\lr{Testbench} و انتخاب ورودی‌ها}
\lr{Testbench} محیطی است که \lr{Device Under Test} یا \lr{DUT} را نمونه‌سازی
می‌کند، ورودی‌ها را اعمال می‌کند و پاسخ را با مرجع مقایسه می‌کند. یک \lr{Test Vector}
می‌تواند شامل ورودی و پاسخ مورد انتظار باشد. در مدار ترتیبی، زمان اعمال ورودی،
\lr{Clock}، \lr{Reset} و وضعیت قبلی نیز بخشی از سناریو هستند
\cite[صفحه ۵۳]{weste}.

برای تابع ترکیبی با \(n\) ورودی دودویی، تعداد حالت‌ها \(2^n\) است، زیرا برای هر
ورودی دو انتخاب مستقل داریم. \lr{Carry} سه‌ورودی ۸ حالت دارد. جمع‌کننده هشت‌بیتی
با دو عملوند هشت‌بیتی و یک \lr{Carry-in}، تعداد
\(2^{8+8+1}=131072\) حالت دارد. با افزایش ورودی‌ها، بررسی کامل به‌سرعت پرهزینه
می‌شود. برای مدار ترتیبی، ترتیب ورودی‌ها و حالت داخلی نیز باید بررسی شوند.

در نمونه آموزشی، برای جمع‌کننده ورودی‌های صفر، بیشینه مقدار، تولید و انتشار
\lr{Carry} و الگوهای مکمل ارزش بررسی دارند. در پردازنده، \lr{Reset}، هر نوع دستور،
شاخه گرفته‌شده و گرفته‌نشده و مرز نشانی‌ها باید طبق تعریف معماری آزموده شوند.
پاسخ مرجع باید از صورت مسئله و محاسبه مستقل به دست آید. کپی همان عبارت اشتباه
از طراحی به مرجع می‌تواند همان خطا را در هر دو حفظ کند.

\lr{Post-synthesis Netlist} را می‌توان با بردارهای آزمون قبلی شبیه‌سازی کرد.
برای شبکه ترانزیستور نیز آزمون عملکرد ممکن است. در مدار ترتیبی، ترتیب رخدادها و
\lr{Race Condition} اهمیت پیدا می‌کند. ابزار \lr{Formal Verification} نیز می‌تواند
تحت فرض‌ها و مدل تعریف‌شده، هم‌ارزی تابع یا خاصیت را بررسی کند.
کیفیت فرض‌ها و پوشش حالت‌ها بر معنای نتیجه اثر می‌گذارد.

\subsection{\lr{DRC}، \lr{LVS}، \lr{ERC} و تحلیل الکتریکی}
\lr{Design Rule Check} یا \lr{DRC} هندسه را با قواعد فناوری مقایسه می‌کند، مانند
حداقل عرض، فاصله و پوشش اتصال. \lr{Layout Versus Schematic} یا \lr{LVS} بررسی
می‌کند که ترانزیستورها و اتصال‌های استخراج‌شده از \lr{Layout} با مدار مرجع سازگار
باشند. \lr{Electrical Rule Check} یا \lr{ERC} برای یافتن مسئله‌های الکتریکی
تعریف‌شده در مجموعه قواعد به کار می‌رود. کتاب به مشکلاتی مانند نویز یا فرسودگی
زودهنگام اشاره می‌کند \cite[صفحه ۵۳]{weste}.

جدول \ref{tab:impl-verification} مثال‌هایی از ارتباط خطا و بررسی مناسب را نشان می‌دهد.
یک مسئله ممکن است در بیش از یک بررسی آشکار شود و توان کشف هر ابزار به مدل و قواعدش وابسته است.
\begin{table}[ht]
\centering
\caption{نمونه ارتباط مسئله طراحی با بررسی لازم.}
\label{tab:impl-verification}
\begin{tabularx}{\textwidth}{rX}
\toprule
مسئله & بررسی مرتبط\\
\midrule
اشتباه در تابع \lr{Carry} & مقایسه رفتاری، جدول کامل یا بررسی هم‌ارزی\\
سیم با عرض کمتر از حد مجاز & \lr{DRC} با قواعد فناوری\\
اتصال خروجی به گره داخلی اشتباه & \lr{LVS} در برابر مدار مرجع\\
اتصال الکتریکی نامناسب تغذیه یا \lr{Body} & \lr{ERC} و بررسی اتصال با قواعد مرتبط\\
رسیدن خروجی پس از مهلت لازم & تحلیل زمان‌بندی با مدل مناسب مدار و سیم\\
مصرف توان بیشتر از بودجه & تحلیل توان با تحریک، فعالیت و مدل معتبر\\
\bottomrule
\end{tabularx}
\end{table}

اگر هندسه قانونی باشد اما دو گره منطقی جابه‌جا شده باشند، ممکن است
\lr{DRC} بدون خطای هندسی پایان یابد و \lr{LVS} اختلاف اتصال را پیدا کند.
اگر اتصال با مرجع برابر باشد اما خود مرجع تابع اشتباه داشته باشد، مقایسه رفتاری
با \lr{Specification} لازم است. برای بررسی سرعت و توان نیز مقاومت و خازن سیم،
اندازه ترانزیستور، ولتاژ و شرایط مدل‌شده اهمیت دارند.

\subsection{تمرین‌های پاسخ‌دار \lr{Verification}}
\textbf{تمرین ۱.} در گیت \lr{NAND2}، خروجی در \lr{Layout} به \(V_{DD}\) اتصال
پیدا کرده است. آزمون عملکرد چه ورودی‌ای را باید حتماً اعمال کند؟
برای \(a=b=1\)، پاسخ مرجع صفر است. خروجی ثابت یک اختلاف را آشکار می‌کند.
سه حالت دیگر خروجی مرجع یک دارند و این خطا را در مدل خروجی ثابت نشان نمی‌دهند.
\lr{LVS} نیز باید اختلاف اتصال با مدار صحیح را بررسی کند.

\textbf{تمرین ۲.} مدار تابع درست را با تأخیر \qty{2}{\nano\second} تولید می‌کند،
ولی مهلت داده‌شده \qty{1.5}{\nano\second} است. نتیجه چیست؟
شرط عملکرد منطقی برآورده شده، اما قید زمان‌بندی \qty{0.5}{\nano\second} نقض شده است.
برای اعلام موفقیت باید هر دو شرط تابع و مهلت برقرار باشند.

\textbf{تمرین ۳.} آیا یک آزمون موفق برنامه \lr{Fibonacci}، همه حالت‌های پردازنده
را پوشش می‌دهد؟ این برنامه تنها توالی مشخصی از دستورها، داده‌ها و شاخه‌ها را طی
می‌کند. برای پوشش نیازهای دیگر باید سناریوهای متناسب با همان نیازها افزوده شود.
مثلاً صحت \lr{Reset}، دستور استفاده‌نشده و حالت مرزی \lr{Branch} به بررسی مربوط خود نیاز دارد.
\clearpage
